% Part: methods
% Chapter: proofs
% Section: proof-by-contradiction

\documentclass[../../../include/open-logic-section]{subfiles}

\begin{document}

\olfileid{mth}{prf}{con}

\olsection{వైరుధ్యం ద్వారా నిరూపణ}

మొదటగా, వైరుధ్యం ద్వారా నిరూపణ ప్రతికూల వాదనలను నిరూపించడానికి ఉపయోగించే నిగమన నమూనా. ఒక వాదన~$p$ \emph{అసత్యం} అని, అంటే~$\lnot p$ అని చూపాలనుకుందాం. ఆశాజనకమైన పద్ధతి: (a) $p$~సత్యమని అనుకోవడం, (b) ఆ పరికల్పన మనకు అసత్యమని తెలిసిన విషయానికి దారితీస్తుందని చూపడం. ఆ ``అసత్యమని తెలిసిన విషయం'' $p$తోనే లేదా మొత్తం వాదనలోని మరొక పరికల్పనతో విరోధించవచ్చు. ఉదాహరణకు ``$q$ అయితే $\lnot p$'' నిరూపణలో $q$~సత్యమని తీసుకొని, దాని నుంచి~$\lnot
p$ను నిరూపించాలి. $\lnot p$ను వైరుధ్యం ద్వారా నిరూపించాలంటే~$q$తో పాటు $p$ను కూడా అనుకోవాలి. $p$ నుంచి $\lnot q$ను నిరూపించగలిగితే, $p$ పరికల్పన మీ మరో పరికల్పన~$q$కు విరుద్ధమైన ఫలితాన్ని ఇస్తుందని చూపినట్లే; ఎందుకంటే $q$, $\lnot q$ రెండూ సత్యం కావు. వైరుధ్యానికి చేరే నిరూపణలో నిర్వచనాలను విప్పడంతో పాటు ఇతర నిగమన నమూనాలనూ ఉపయోగించాలి. ఒక ఉదాహరణ చూద్దాం.

\begin{prop}
  $A \subseteq B$, $B = \emptyset$ అయితే $A$లో \tetoken{మూలకాలు}{!!{element}s} ఉండవు.
\end{prop}

\begin{proof}
  $A \subseteq B$, $B = \emptyset$ అనుకుందాం. $A$లో \tetoken{మూలకాలు}{!!{element}s} లేవని చూపాలి.
  \begin{quote}
    ఇది సోపాధిక వాదన కాబట్టి పూర్వపక్షాన్ని తీసుకొని పర్యవసానాన్ని నిరూపించాలి. పర్యవసానం: $A$లో \tetoken{మూలకాలు}{!!{element}s} లేవు. దీన్ని మరింత స్పష్టంగా, $x \in A$ అయ్యేదేదీ లేదని చెప్పవచ్చు.
  \end{quote}
  $A$లో \tetoken{మూలకాలు}{!!{element}s} లేవు అప్పుడూ అప్పుడే $x \in A$ అయ్యే~$x$ ఒకటి ఉందనేది అసత్యం.
  \begin{quote}
    నిరూపించాల్సింది నిజానికి ప్రతికూల వాదన~$\lnot p$: $x \in A$ అయ్యేదేదీ లేదు. వైరుధ్యం ద్వారా నిరూపణలో సంబంధిత సానుకూల వాదన~$p$, అంటే $x \in A$ అయ్యేదొకటి ఉందని తీసుకొని దాని నుంచి వైరుధ్యాన్ని నిరూపించాలి. దీన్ని సూచించడానికి ``వైరుధ్యం కోసం అనుకుందాం'' లేదా ``కాదనుకుందాం'' అని రాసి, తరువాత పరికల్పన~$p$ను చెబుతాం.
  \end{quote}
  కాదనుకుందాం: $x \in A$ అయ్యేది ఒకటి ఉంది.
  \begin{quote}
    వైరుధ్యానికి చేరేందుకు ఇప్పుడు ఈ కొత్త పరికల్పనను వాడతాం. మరో రెండు పరికల్పనలు: $A \subseteq B$, $B = \emptyset$. మొదటిదాని నుంచి $x \in B$ వస్తుంది:
  \end{quote}
  $A \subseteq B$ కనుక $x \in B$.
  \begin{quote}
    అయితే $B = \emptyset$ కనుక $B$లోని ప్రతి \tetoken{మూలకం}{!!{element}} (ఉదా., $x$) కూడా~$\emptyset$లో \tetoken{ఒక మూలకం}{!!a{element}} కావాలి.
  \end{quote}
  $B = \emptyset$ కనుక $x \in \emptyset$. ఇది వైరుధ్యం; ఎందుకంటే నిర్వచనం ప్రకారం $\emptyset$లో \tetoken{మూలకాలు}{!!{element}s} లేవు.
  \begin{quote}
    నిరూపణ ఇక్కడే పూర్తవుతుంది: తీసుకున్న పరికల్పనల నుంచి అవసరమైన వైరుధ్యం వచ్చింది; కాబట్టి అవన్నీ ఒకేసారి సత్యం కాలేవు. మొదటి రెండు పరికల్పనలను ($A \subseteq B$, $B = \emptyset$) ప్రశ్నించడం లేదు కనుక, చివరగా ప్రవేశపెట్టిన పరికల్పన ($x \in A$ అయ్యేది ఉంది) అసత్యం కావాలి. మరింత స్పష్టత కోసం దాన్ని విడిగా చెప్పవచ్చు.
  \end{quote}
  అందువల్ల $x \in A$ అయ్యేదొకటి ఉందన్న పరికల్పన అసత్యం; కనుక వైరుధ్యం ద్వారా నిరూపణ ప్రకారం $A$లో \tetoken{మూలకాలు}{!!{element}s} లేవు.
\end{proof}

సాంప్రదాయిక తర్కంలో ప్రతి సానుకూల వాదన ఒక ప్రతికూల వాదనతో సమానార్థకం: $p$ అప్పుడూ అప్పుడే $\lnot\lnot p$. అందువల్ల వైరుధ్యం ద్వారా నిరూపణను సానుకూల వాదనలను ``పరోక్షంగా'' స్థాపించడానికి కూడా వాడవచ్చు. $p$ను నిరూపించాలంటే దాన్ని $\lnot\lnot p$ అనే ప్రతికూల వాదనగా చదవాలి. $\lnot p$ నుంచి వైరుధ్యాన్ని నిరూపించగలిగితే, వైరుధ్యం ద్వారా $\lnot\lnot p$ను, సాంప్రదాయిక ద్వినిషేధ నిర్మూలనతో~$p$ను స్థాపించినట్లే.
\sourcecorrection{OLTEMTHPRFCON-001}{మూలం $p$ మరియు $\lnot\lnot p$ సమానార్థకతను ఎలాంటి పరిధి లేకుండా ``స్వయంగా'' అంది. సాంప్రదాయిక తర్కంలో ఈ పరోక్ష సానుకూల నిరూపణ చెల్లుతుంది; అంతఃప్రజ్ఞావాద తర్కంలో సాధారణ ద్వినిషేధ నిర్మూలన చెల్లదు. ఆ పరిధిని స్పష్టంచేశాం.}

గత ఉదాహరణలో ప్రతికూల వాదనను, అంటే $A$లో \tetoken{మూలకాలు}{!!{element}s} లేవని, నిరూపించాలనుకున్నాం. అందుకే వైరుధ్యం ద్వారా నిరూపణ కోసం తీసుకున్న పరికల్పన ($x \in A$ అయ్యేది ఉంది) సానుకూల వాదన. అది తర్కించడానికి ఊహాత్మక $x \in A$ను ఇచ్చింది; దాని గురించి సాగి $x \in \emptyset$కు చేరుకున్నాం.

సానుకూల వాదనను పరోక్షంగా నిరూపించేటప్పుడు వైరుధ్యం కోసం తీసుకునే పరికల్పన ప్రతికూలంగా ఉంటుంది. అయితే సానుకూల వాదనను ప్రతికూల వాదనగా, ప్రతికూల వాదనను సానుకూల వాదనగా తరచుగా సులభంగా తిరిగి చెప్పవచ్చు. గత నిరూపణలో ``$A$లో \tetoken{మూలకాలు}{!!{element}s} లేవు'' అనే ప్రతికూల పర్యవసానం బదులు ``$A = \emptyset$''ను నిరూపించినా నిరూపణ సారాంశం అదే. ($=$ నిర్వచనం ప్రకారం ``$A = \emptyset$'' అనేది సార్వత్రిక వాదన; దాన్ని విప్పితే ``$A$లోని ప్రతి \tetoken{మూలకం}{!!{element}},~$\emptyset$లో \tetoken{ఒక మూలకం}{!!{element}}; అలాగే వ్యతిరేక దిశలోనూ'' అవుతుంది.) అయితే అది ``$x \in A$ అయ్యేదేదీ లేదు'' అనే ప్రతికూల వాదనకు సమానార్థకమని సులభంగా చూడవచ్చు.

కాబట్టి కొన్నిసార్లు $p$ను నేరుగా నిరూపించడంకన్నా $\lnot p$ను పరికల్పనగా తీసుకోవడం సులభం. నేర నిరూపణ అంతే సరళంగా లేదా ఇంకా సరళంగా ఉన్నా (తరువాతి ఉదాహరణలలాగా) కొందరు పరోక్ష మార్గాన్నే ఇష్టపడతారు. ద్వినిషేధం గందరగోళంగా ఉంటే, ఒక వాదనను వైరుధ్యం ద్వారా నిరూపించడం అంటే దాని \emph{వ్యతిరేక} వాదన నుంచి వైరుధ్యాన్ని నిరూపించడం అని భావించండి. కాబట్టి $\lnot
p$కు వైరుధ్యం ద్వారా నిరూపణ అంటే~$p$ పరికల్పన నుంచి వైరుధ్యం;~$p$కు వైరుధ్యం ద్వారా నిరూపణ అంటే~$\lnot p$ నుంచి వైరుధ్యం.

\begin{prop}
$A \subseteq A \cup B$.
\end{prop}

\begin{proof}
  $A \subseteq A \cup B$ అని చూపాలి.
  \begin{quote}
    పైకి చూస్తే ఇది సానుకూల వాదన: ప్రతి $x \in A$, $A \cup B$లోనూ ఉంది. దీని నిషేధం: ఏదో $x \in A$కు $\notin A \cup B$. కాబట్టి ఈ నిషేధాన్ని తీసుకొని అది వైరుధ్యానికి దారితీస్తుందని చూపి వాదనను పరోక్షంగా నిరూపించవచ్చు.
  \end{quote}
  కాదనుకుందాం; అంటే $A \nsubseteq A \cup B$.
  \begin{quote}
    $A \subseteq A \cup B$ నిర్వచనం మనకు తెలుసు: ప్రతి $x \in A$ కూడా $\in A \cup B$. $A \nsubseteq A \cup B$ అర్థం తెలుసుకోవడానికి విప్పిన నిర్వచనంపై ప్రాథమిక తార్కిక మార్పు చేయాలి: ప్రతి $x \in A$ కూడా $\in
    A \cup B$ అన్నది అసత్యం అప్పుడూ అప్పుడే \emph{ఏదో}~$x \in A$కు $\notin A \cup B$. (ఇక్కడ చాలా జాగ్రత్త అవసరం: ``అన్ని $A$s, $B$s'' వంటి వాదన నిషేధాన్ని తప్పుగా విశ్లేషించడం వల్ల విద్యార్థుల వైరుధ్య నిరూపణలు తరచుగా విఫలమవుతాయి.) మరో మాటలో $A \nsubseteq A \cup B$ అప్పుడూ అప్పుడే $x \in A$, $x \notin A \cup B$ అయ్యే $x$ ఒకటి ఉంది. ఆ తరువాత సులభం.
    \sourcecorrection{OLTEMTHPRFCON-002}{మూలంలో ``ఏదో $x\in A$కు $\notin C$'' అని ఉంది; ఈ ఉదాహరణ లక్ష్య సమితి $A\cup B$, ఇక్కడ $C$ నిర్వచించలేదు. $\notin A\cup B$గా సరిచేశాం.}
  \end{quote}
  కాబట్టి $x \in A$, $x \notin A \cup B$ అయ్యే వస్తువు ఒకటి ఉంది. $\cup$ నిర్వచనం ప్రకారం $x \in A \cup B$ అప్పుడూ అప్పుడే $x \in A$ లేదా $x \in
  B$. $x \in A$ కనుక $x \in A \cup B$. ఇది $x \notin A \cup B$ పరికల్పనకు విరుద్ధం.
\end{proof}

\begin{prob}
$A \cap B \subseteq A$ అని \emph{పరోక్షంగా} నిరూపించండి.
\end{prob}

\begin{prop}
$A \subseteq B$, $B \subseteq C$ అయితే $A \subseteq C$.
\end{prop}

\begin{proof}
  $A \subseteq B$, $B \subseteq C$ అనుకుందాం. $A
  \subseteq C$ అని చూపాలి.
  \begin{quote}
    పరోక్షంగా సాగుదాం: నిరూపించాల్సిన దాని నిషేధాన్ని తీసుకుందాం.
  \end{quote}
  కాదనుకుందాం; అంటే $A \nsubseteq C$.
  \begin{quote}
    ముందులాగే $A \nsubseteq C$ అంటే ప్రతి $x \in A$ కూడా $\in C$ అన్నది అసత్యం; అంటే ఏదో $x \in A$కు $\notin C$. సాధనతో ఇలాంటి నిషేధాలను విప్పడానికి ఎక్కువ ఆలోచించనవసరం ఉండదు.
  \end{quote}
  మరో మాటలో $x \in A$, $x \notin C$ అయ్యే~$x$ ఒకటి ఉంది.
  \begin{quote}
    ఇప్పుడు దీన్ని వైరుధ్యానికి చేరడానికి వాడవచ్చు. అందుకు మిగిలిన రెండు పరికల్పనలనూ ఉపయోగించాలి.
  \end{quote}
  $A \subseteq B$ కనుక $x \in B$. $B \subseteq C$ కనుక $x \in
  C$. ఇది $x \notin C$కు విరుద్ధం.
\end{proof}

\begin{prop}
$A \cup B = A \cap B$ అయితే $A = B$.
\end{prop}

\begin{proof}
  $A \cup B = A \cap B$ అనుకుందాం. $A = B$ అని చూపాలి.
  \begin{quote}
    ప్రారంభం ఇప్పుడు సాధారణమే:
  \end{quote}
  వైరుధ్యం కోసం $A \neq B$ అనుకుందాం.
  \begin{quote}
    వైరుధ్యం కోసం తీసుకున్న పరికల్పన $A \neq
    B$. $A = B$ అప్పుడూ అప్పుడే $A \subseteq B$ మరియు $B \subseteq A$ కనుక, $A \neq B$ అప్పుడూ అప్పుడే $A \nsubseteq B$ \emph{లేదా} $B \nsubseteq
    A$. (నిషేధాలను మార్చేటప్పుడు జాగ్రత్త ఎంత ముఖ్యమో గమనించండి!{}) ఈ వికల్పం నుంచి వైరుధ్యాన్ని చూపడానికి సందర్భాలవారీ నిరూపణ చేసి, ప్రతి సందర్భంలో వైరుధ్యం వస్తుందని చూపుతాం.
  \end{quote}
  $A \neq B$ అప్పుడూ అప్పుడే $A \nsubseteq B$ లేదా $B \nsubseteq A$. సందర్భాలను వేరుచూద్దాం.
  \begin{quote}
    మొదటి సందర్భంలో $A \nsubseteq B$ అనుకుందాం; అంటే ఏదో $x$కు $x \in A$ కానీ $\notin B$. $A \cap B$ అనేది $A$, $B$ రెండింటిలోనూ ఉన్న \tetoken{మూలకాల}{!!{element}s} సమితి; కనుక ఒకదానిలో లేకపోతే ఛేదనలో ఉండదు. $A \cup B$ అనేది $A$తో పాటు $B$; రెండింటిలో ఏదో ఒకదానిలో ఉన్నది సమ్మేళనంలో ఉంటుంది. కాబట్టి $x \in A \cup B$ కానీ $x \notin A \cap
    B$; అందువల్ల $A \cap B \neq A \cup B$.
  \end{quote}
  
  సందర్భం 1: $A \nsubseteq B$. అప్పుడు ఏదో $x$కు $x \in A$ కానీ $x \notin
  B$. $x \notin B$ కనుక $x \notin A \cap B$. $x \in A$ కనుక $x \in A \cup B$. కాబట్టి $A \cap B \neq A \cup B$; ఇది $A \cap B = A \cup B$ పరికల్పనకు విరుద్ధం.

  సందర్భం 2: $B \nsubseteq A$. అప్పుడు ఏదో $y$కు $y \in B$ కానీ $y \notin
  A$. ముందులాగే $y \in A \cup B$ కానీ $y \notin A \cap B$; కాబట్టి $A \cap B \neq A \cup B$, మళ్లీ $A \cap B = A \cup
  B$కు విరుద్ధం.
\end{proof}

\end{document}
